% ======================================================================
% EXPANSAO — Triangulacao Anti-Circularidade (SPEC-008)
% Framework para dominios sem ground truth externo
% 15 referencias com DOI | TDD: 14/14 GREEN
% ======================================================================

\section{Triangulacao Anti-Circularidade: Validacao sem Ground Truth Externo}

\subsection{O Problema da Circularidade Auto-Avaliativa}

O dialogo com o revisor senior simulado Romulo identificou a questao mais
profunda da validacao cientifica em IA: \textbf{e quando nao existe um Project
Euler para o dominio?} Quando o unico ground truth disponivel e interno ao
proprio corpus — como um extrator de padroes em textos juridicos
especializados — como sair da circularidade em que o sistema valida a si mesmo?

Este problema e formalizado pela Lei de Goodhart\footnote{Goodhart, C.A.E.
\textit{Problems of Monetary Management: The UK Experience}. Papers in Monetary
Economics, Reserve Bank of Australia, 1975.} e sua generalizacao por
Strathern\footnote{Strathern, M. \textit{'Improving ratings': audit in the
British University system}. European Review, 5(3):305--321, 1997. DOI:
10.1002/(SICI)1234-981X(199707)5:3<305::AID-EURO184>3.0.CO;2-4.}: "quando uma
medida se torna um alvo, ela deixa de ser uma boa medida". Ji et
al.\footnote{Ji, J. et al. \textit{AI Alignment: A Comprehensive Survey}.
arXiv:2310.19852, 2023. DOI: 10.48550/arxiv.2310.19852.} catalogam este
fenomeno como ``reward hacking'' — o sistema encontra atalhos para maximizar
a metrica sem desenvolver a capacidade subjacente.

No contexto do CORA-Eval, a circularidade se manifesta quando os verificadores
Cora-Debate (V1--V7) sao usados simultaneamente como mecanismo de verificacao
das tarefas e como componentes do ecossistema sendo avaliado. Para as dimensoes
D1 e D7, esta circularidade e mitigada pela validacao externa (Project Euler e
Rosalind). Para as 8 dimensoes restantes — e para qualquer dominio sem ground
truth externo — permanece a questao: como quebrar a circularidade?

\subsection{Fundamentacao Teorica: Triangulacao de Denzin-Flick}

A resposta metodologica fundamenta-se na triangulacao, proposta por
Denzin\footnote{Denzin, N.K. \textit{The Research Act: A Theoretical
Introduction to Sociological Methods}. McGraw-Hill, 1978. ISBN: 978-0070168904.}
e atualizada por Flick\footnote{Flick, U. \textit{Triangulation in Data
Collection}. In: The SAGE Handbook of Qualitative Data Collection, 2018. DOI:
10.4135/9781526416070.}: multiplos metodos, multiplas fontes e multiplos
investigadores convergem para reduzir o vies de cada um isoladamente.

Adaptamos este framework para o problema da validacao circular em IA,
substituindo ``investigadores'' por ``fontes de ground truth'' e ``metodos''
por ``camadas de verificacao com independencia progressiva''. O resultado e o
\textbf{Framework de Triangulacao Anti-Circularidade} (SPEC-008), com tres
camadas de validacao.

\subsection{As Tres Camadas}

\subsubsection{Camada 1 — Split Temporal Cego}

Seja $\mathcal{D} = \{(x_i, t_i)\}_{i=1}^{N}$ um corpus com timestamps. O
split temporal cego particiona:

\begin{equation}
\mathcal{D}_{\text{treino}} = \{x_i : t_i \leq T_{\text{cutoff}}\}, \qquad
\mathcal{D}_{\text{teste}} = \{x_i : t_i > T_{\text{cutoff}}\}
\end{equation}

O modelo de extracao $\mathcal{M}$ e treinado exclusivamente em
$\mathcal{D}_{\text{treino}}$ e avaliado em $\mathcal{D}_{\text{teste}}$, que
contem dados inexistentes no momento do treinamento. Se o padrao extraido do
periodo 2015--2023 \textbf{prediz} corretamente o que aparece em 2024--2025,
ha evidencia de que captura algo real.

Bergmeir e Benitez\footnote{Bergmeir, C., Benitez, J.M. \textit{On the use of
cross-validation for time series predictor evaluation}. Information Sciences,
191:192--213, 2012. DOI: 10.1016/j.ins.2011.12.028.} demonstraram que validacao
cruzada convencional (K-fold aleatorio) produz estimativas excessivamente
otimistas para dados com dependencia temporal — o futuro vaza para o passado.
Cerqueira et al.\footnote{Cerqueira, V., Torgo, L., Mozetic, I. \textit{Evaluating
time series forecasting models}. Machine Learning, 109:1997--2028, 2020. DOI:
10.1007/s10994-020-05910-7.} confirmaram que CV com janela deslizante
preservando ordem temporal supera K-fold aleatorio em 62 series temporais reais.

\textbf{Nivel de independencia:} BAIXO (mesmo corpus, mesma fonte, separacao
temporal). Nao elimina a circularidade, mas a enfraquece ao introduzir o tempo
como barreira.

\subsubsection{Camada 2 — Perturbacao Adversaria (Falsificabilidade)}

A Camada 2 operacionaliza o criterio de falsificabilidade de
Popper\footnote{Popper, K. \textit{The Logic of Scientific Discovery}.
Routledge, 1959. DOI: 10.4324/9780203994627.} no contexto de IA: em vez de
tentar provar que o padrao e correto, tentamos \textit{refutar} o padrao por
perturbacao.

Para cada padrao $P$ extraido do corpus original, criamos $k=4$ copias
perturbadas aplicando transformacoes $T_j$ que destroem propriedades
especificas:

\begin{itemize}
\item $T_1$: embaralhar ordem de paragrafos (destroi coerencia discursiva)
\item $T_2$: substituir entidades nomeadas por \texttt{[ENT\_\#\#\#]}
      (destroi semantica referencial)
\item $T_3$: inverter ordem cronologica (destroi estrutura temporal)
\item $T_4$: substituir 30\% de tokens por sinonimos aleatorios (destroi
      precisao terminologica)
\end{itemize}

Se $P$ \textbf{desaparece} sob $T_j$, a propriedade $j$ e necessaria para $P$
— o padrao nao e espurio. Se $P$ \textbf{persiste} identico sob todas as
perturbacoes, $P$ e insensivel a estrutura do corpus e provavelmente captura
um vies de distribuicao de tokens (artefato, nao conhecimento).

Ribeiro et al.\footnote{Ribeiro, M.T., Wu, T., Guestrin, C., Singh, S.
\textit{Beyond Accuracy: Behavioral Testing of NLP Models with CheckList}. ACL,
2020. DOI: 10.18653/v1/2020.acl-main.442.} formalizaram esta abordagem no
CheckList, com tres tipos de teste: Minimum Functionality Test (MFT),
Invariance Test (INV) e Directional Expectation Test (DIR). Gardner et
al.\footnote{Gardner, M. et al. \textit{Evaluating Models' Local Decision
Boundaries via Contrast Sets}. EMNLP Findings, 2020. DOI:
10.18653/v1/2020.findings-emnlp.117.} propoem contrast sets — exemplos que
diferem minimamente mas exigem saidas diferentes.

\textbf{Nivel de independencia:} MEDIO (mesmo corpus, condicoes modificadas).

\subsubsection{Camada 3 — Anotacao Humana em Amostra Minima}

A Camada 3 introduz um especialista humano do dominio como fonte de ground
truth externa ao sistema. Dado que anotacao humana e cara, utilizamos
\textbf{active learning}\footnote{Settles, B. \textit{Active Learning Literature
Survey}. Computer Sciences Technical Report 1648, University of
Wisconsin-Madison, 2009.}: o sistema seleciona para anotacao os $m=30$
documentos onde o modelo tem maior incerteza (entropia maxima), ha divergencia
entre C1 e C2, ou o documento e representativo de um cluster ainda nao anotado.

Shen et al.\footnote{Shen, Y., Yun, H., Lipton, Z.C., Kronrod, Y., Anandkumar,
A. \textit{Deep Active Learning for Named Entity Recognition}. ICLR, 2018. DOI:
10.48550/arxiv.1707.05928.} demonstram que estrategias hibridas (incerteza +
representatividade) superam uncertainty sampling puro, reduzindo o custo de
rotulagem em ate 80\% comparado a amostragem aleatoria.

O protocolo de anotacao consiste em 3 perguntas binarias por padrao por
documento: (1) O padrao esta presente? (2) Captura algo real ou e artefato?
(3) O especialista teria identificado sem o sistema? A concordancia e medida
como proporcao de respostas positivas em (1), com intervalo de confianca de
95\% via Clopper-Pearson exato (binomial, nao aproximacao normal de Wald).

\textbf{Nivel de independencia:} ALTO (fonte externa: julgamento humano).

\subsection{Matriz de Decisao Integrada}

\begin{table}[H]\centering\footnotesize
\caption{Matriz de decisao da triangulacao anti-circularidade}
\label{tab:triangulacao}
\begin{tabular}{p{0.05\textwidth}p{0.10\textwidth}p{0.12\textwidth}p{0.10\textwidth}p{0.30\textwidth}p{0.15\textwidth}}
\toprule
\textbf{Cen.} & \textbf{C1 (Temp)} & \textbf{C2 (Pert)} & \textbf{C3 (Hum)} & \textbf{Interpretacao} & \textbf{Acao} \\
\midrule
A & Passa & Passa & Passa & Evidencia forte de robustez & Publicar com confianca moderada \\
B & Passa & Falha & Passa & Padrao sensivel a estrutura & Investigar perturbacao que quebrou \\
C & Falha & Passa & Passa & Overfitting temporal & Verificar domain shift \\
D & Passa & Passa & Falha & Vies do especialista ou padrao irrelevante & Discutir, expandir amostra \\
E & Falha & Falha & Passa & Padrao fragil & Refinar extracao, nao publicar \\
F & — & — & Falha & Discordancia humana & Expandir amostra para 60 \\
\bottomrule
\end{tabular}
\end{table}

\subsection{Niveis de Independencia}

\begin{table}[H]\centering\footnotesize
\caption{Niveis de independencia e confianca por fonte de validacao}
\label{tab:independencia}
\begin{tabular}{p{0.30\textwidth}p{0.15\textwidth}p{0.20\textwidth}p{0.20\textwidth}}
\toprule
\textbf{Fonte de Validacao} & \textbf{Independencia} & \textbf{Confianca Max.} & \textbf{Custo} \\
\midrule
Cora-Debate auto-verif. & Nula (circular) & — & Zero \\
Split temporal cego (C1) & Baixa & Moderada & Baixo \\
Perturbacao adversaria (C2) & Media & Moderada-Alta & Medio \\
Anotacao humana 30 amostras (C3) & Alta & Alta & 2h humanas \\
Validacao externa (Project Euler) & Maxima & Maxima & Depende de existencia \\
Replicacao por terceiros & Maxima & Maxima & Meses \\
\bottomrule
\end{tabular}
\end{table}

\subsection{Integracao com CORA-Eval e TDD (SPEC-008)}

A triangulacao anti-circularidade foi formalizada como a especificacao
SPEC-008 do ecossistema, com 9 criterios de teste (CT-8.1 a CT-8.9)
implementados em \texttt{test\_anticircularidade.py}. A suite TDD obteve
14/14 GREEN apos 3 correcoes no ciclo RED$\to$GREEN$\to$REFACTOR.

A integracao com o CORA-Eval opera como camada adicional de verificacao
que precede os verificadores Cora-Debate:

\begin{verbatim}
Triangulacao Anti-Circularidade (C1+C2+C3)
    |
    v
Cora-Debate V1-V7 (verificacao simbolica)
    |
    v
Qualis A1 Scoring
\end{verbatim}

A triangulacao estabelece o \textbf{nivel de independencia} da validacao.
O Cora-Debate opera dentro desse nivel, com a ressalva de que scores
gerados em dominios sem validacao externa recebem a penalidade
\texttt{[auto-reportado]} conforme INTEGRIDADE.md R-I8.

\subsection{Limitacoes do Framework}

\begin{enumerate}[label=(\roman*)]
\item A Camada 3 depende de disponibilidade de especialista humano — custo
      real de 2 horas, nao simulavel.
\item O split temporal requer timestamps confiaveis no corpus — nem todo
      dominio possui ordenacao temporal.
\item As perturbacoes T1--T4 sao heuristicas — podem nao cobrir todos os
      tipos de fragilidade.
\item A matriz de decisao A-F e qualitativa — os thresholds sao heuristicos
      e podem requerer calibracao por dominio.
\item A triangulacao \textbf{reduz}, mas nao \textbf{elimina}, a
      circularidade. Nenhuma das tres camadas substitui uma validacao
      externa independente.
\end{enumerate}

\subsection{Referencias do Framework}

\begin{enumerate}[label={[\arabic*]}]
\item Goodhart, C.A.E. (1975). Problems of Monetary Management. Reserve Bank
      of Australia.
\item Strathern, M. (1997). `Improving ratings'. European Review, 5(3).
      DOI: 10.1002/(SICI)1234-981X(199707)5:3.
\item Ji, J. et al. (2023). AI Alignment: A Comprehensive Survey.
      arXiv:2310.19852. DOI: 10.48550/arxiv.2310.19852.
\item Denzin, N.K. (1978). The Research Act. McGraw-Hill.
      ISBN: 978-0070168904.
\item Flick, U. (2018). Triangulation in Data Collection. SAGE Handbook.
      DOI: 10.4135/9781526416070.
\item Bergmeir, C., Benitez, J.M. (2012). Cross-validation for time series.
      Information Sciences, 191. DOI: 10.1016/j.ins.2011.12.028.
\item Cerqueira, V. et al. (2020). Evaluating time series forecasting models.
      Machine Learning, 109. DOI: 10.1007/s10994-020-05910-7.
\item Arlot, S., Celisse, A. (2010). Cross-validation procedures. Statistics
      Surveys, 4. DOI: 10.1214/09-SS054.
\item Ribeiro, M.T. et al. (2020). CheckList. ACL.
      DOI: 10.18653/v1/2020.acl-main.442.
\item Gardner, M. et al. (2020). Contrast Sets. EMNLP Findings.
      DOI: 10.18653/v1/2020.findings-emnlp.117.
\item Popper, K. (1959). The Logic of Scientific Discovery. Routledge.
      DOI: 10.4324/9780203994627.
\item Settles, B. (2009). Active Learning Literature Survey. UW-Madison
      CS TR 1648.
\item Shen, Y. et al. (2017). Deep Active Learning for NER. ICLR 2018.
      DOI: 10.48550/arxiv.1707.05928.
\item Hendrycks, D., Gimpel, K. (2017). Out-of-Distribution Detection.
      ICLR. DOI: 10.48550/arxiv.1610.02136.
\item Goodfellow, I. et al. (2015). Adversarial Examples. ICLR.
      DOI: 10.48550/arxiv.1412.6572.
\end{enumerate}
